\chapter{Introduction}
\label{ch:intro}

\section{The problem}

Two articles can report the same event and leave a reader in entirely different states. One
informs; the other provokes. Readers notice the difference and cannot quantify it, and the
platforms that distribute both have no objective measure of which they are amplifying. The
difference is real, consequential and unmeasured.

Existing approaches measure the text. Sentiment analysis counts affective vocabulary;
readability scores count syllables and clause length; stance detection classifies a position.
None of them measures what the content is predicted to \emph{do} to the person reading it.

This project takes a different route. Rather than characterising the words, it estimates the
response: a released neural encoder predicts cortical activity from the content, and the
measurement is taken on that prediction. The instrument is a thermometer for one specific
property of media, and the thesis is an account of building it, running it, and stating
carefully what its readings support.

\section{The amendment, stated at the outset}
\label{sec:amendment}

Three commitments in the original proposal could not be met as written. They are set out here
rather than disclosed in the results chapter, because a reader who meets them for the first
time beside a result will reasonably suspect they were discovered to explain it. Each was
identified before the measurement was taken, and each is documented in the amendment approved
by the supervisor.

\textbf{First, the observable is cortical.} The proposal framed the index around subcortical
affective structures. The released checkpoint predicts cortical surface activity only, over
$20{,}484$ vertices, and contains no subcortical output. The index is therefore a
\emph{cortical proxy}: a contrast between two cortical ROI unions, one associated with
affective salience and one with deliberative control. No claim about subcortical structures
appears anywhere in this document.

\textbf{Second, the index is a difference, not a ratio.} Equation~(4) of the proposal defines
the index as $A_{\mathrm{aff}}/(A_{\mathrm{del}} + \delta)$. The encoder emits standardised
activation, so an ROI mean is negative about as often as positive and the ratio sign-flips or
diverges. On the completed corpus it is undefined for $17.25\%$ of items. The measurement uses
the signed difference $A_{\mathrm{aff}} - A_{\mathrm{del}}$, which is defined everywhere and
is the form the physics requires.

\textbf{Third, the corpus is text-only.} The proposal specified 1{,}500 multimodal items. The
delivered corpus is 400 text items in four balanced categories, and Research Question~IV,
which required audio and video, is out of scope.

\section{Research questions}

\begin{description}
\item[RQ I] Can predicted cortical activation serve as a quantitative signature of
emotionally manipulative media content, and how does the resulting index perform as a
classifier against human labels?

\item[RQ II] How is the index distributed across content categories, and do the categories
separate?

\item[RQ III] What phase structure does a mean-field opinion model exhibit under a field
derived from this observable, and what does the measured spread imply for the coupling?

\item[RQ IV] How do text, audio and video contribute to the index? \emph{Out of scope under
the amendment.}
\end{description}

\section{Objectives}

\begin{enumerate}[label=(\roman*)]
\item Build a labelled corpus of media items in balanced categories.
\item Run the encoder over the corpus and compute the index for every item.
\item Characterise the distribution of the index.
\item Calibrate the coupling between the index and the model's field.
\item Derive the phase structure of the opinion model.
\item Validate the index as a classifier against human labels.
\item Deliver the measurement apparatus as working open-source code.
\end{enumerate}

Objective (vii) is the primary deliverable. The proposal states this explicitly and it governs
how the work should be assessed: the apparatus is the promise, and the empirical result is
what the apparatus returned when it was run.

\section{Contribution}

Three, developed across the chapters that follow.

A \textbf{bound} on the coupling between any content observable and an opinion-dynamics field,
requiring only the observable's measured spread and the model's own spinodal, and evaluable
before data are collected (Chapter~3).

An \textbf{instrument}, delivered as open-source code, together with the documented finding
that the released checkpoint is cortical-only and that its agreement with real cortex is
contested (Chapter~4).

A \textbf{measurement} over 400 items, reported with the design's sensitivity attached and its
central confound stated before its interpretation (Chapters~5 and~6).

\section{Structure of this document}

Chapter~2 reviews the literature on opinion dynamics as statistical physics and on predictive
neural encoding, and identifies the gap this work addresses. Chapter~3 derives the
theoretical framework and the bound. Chapter~4 describes the instrument, the corpus and the
analysis. Chapter~5 reports the results. Chapter~6 discusses what they support. Chapter~7
concludes.
